\documentclass[12pt, a4paper]{article}

\usepackage[utf8]{inputenc}
\usepackage[T1]{fontenc}

\DeclareUnicodeCharacter{00B0}{\ensuremath{^\circ}}
\usepackage{amsmath, amssymb}
\usepackage{graphicx}
\usepackage{xcolor}
\usepackage{booktabs}
\usepackage{natbib}
\usepackage{hyperref}
\usepackage[margin=2.5cm]{geometry}
\usepackage{lineno}
\usepackage{setspace}
\usepackage{siunitx}
\usepackage{cleveref}
\usepackage{placeins}
\usepackage{rotating}

\doublespacing
\linenumbers


\newcommand{\todo}[1]{\textcolor{red}{\textbf{[TODO: #1]}}}
\newcommand{\placeholder}[1]{\textcolor{gray}{\textit{#1}}}
\newcommand{\Sv}{\,\text{Sv}}    % 1 Sv = 10^6 m^3/s
\newcommand{\psu}{\,\text{psu}}
\newcommand{\CO}{CO\ensuremath{_2}}
\newcommand{\dint}{\mathrm{d}}


\begin{document}

\title{Coherence of Resilience Measures for the Atlantic Meridional
       Overturning Circulation Across a Hierarchy of Models}

\author{Andreas Morr$^{1,2}$%
\thanks{Corresponding author: \texttt{andreas.morr@tum.de}}
\and
Reyk Börner$^{3}$
\and
Sjoerd Terpstra$^{3}$
\and
Christian Kuehn$^{1}$}


\date{\today}
\maketitle

\begin{abstract}
The Atlantic Meridional Overturning Circulation (AMOC) is a critical component of the global climate system, and its potential collapse one of the most consequential climate tipping events. Quantifying AMOC resilience, that is, its capacity to maintain the active circulation state under external perturbations, is therefore of high scientific and societal importance. Resilience is a multifaceted concept, encompassing many notions of transient system response. Whether these distinct measures yield consistent, and therefore trustworthy, signals across models of different complexity has not been systematically examined. Here, we address this question by analysing four models that span a hierarchy in physical complexity: (i) a 3-box salinity model, (ii) a two-dimensional latitude-depth Boussinesq model of the overturning circulation, (iii) the Earth system model of intermediate complexity CLIMBER-X, and (iv) the general circulation model PlaSim-LSG. Across the hierarchy, we compute model-specific implementations of four complementary resilience concepts: local resilience, convergence time, basin volume, and minimal critical shock. Our results suggest that not all measures are coherent across the model hierarchy, as exhibited by incompatible trends under increasing \CO{} forcing in each model. This points to deficiencies in model compatibility and challenges in applying these concepts in real-world systems. Several measures indicate an unexpected increase in system resilience under climate change. These findings have direct implications for the detectability of an imminent AMOC collapse from observational or model-based early-warning signals.

\end{abstract}

\textbf{Keywords:} Atlantic Meridional Overturning Circulation, tipping points, resilience, model hierarchy

\newpage

\section{Introduction}
\label{sec:intro}

The Atlantic Meridional Overturning Circulation is one of the dominant features of the Earth's oceans, transporting large amounts of warm, saline surface waters northward across the Atlantic and returning cold, deep North Atlantic water southward at depth \citep{Talley2003DataBasedMeridionalOverturning, Johns2023AMOCHeatTransport}. This large-scale transport pattern underpins the relatively mild climate of Western Europe and plays a central role in the global carbon and heat cycles \citep{Ganachaud2000ImprovedEstimatesGlobal, Jackson2015AMOCImpact, Schaumann2025AMOCCarbonCosts}. Palaeoclimate records document repeated rapid transitions between an active (AMOC-on) and a weakened or collapsed (AMOC-off) state during the Pleistocene, associated with Dansgaard–Oeschger events \citep{Dansgaard1984DO, Clark2002AMOCAbruptClimate, Menviel2020DOIceOscillation}. Under continued anthropogenic forcing, rising sea surface temperatures, freshening of the North Atlantic from increased precipitation, and Greenland ice-sheet melt may push the AMOC towards a similar transition \citep{Yang2016AMOCFreshwater, Lohmann2021RiskTippingOverturning, ArmstrongMcKay2022TippingPoints, VanWesten2024PhysicsBasedEWS}.

More concretely, in background climate configurations where at least two stable equilibria coexist (the on and off states), the system can be permanently tipped from one to the other by crossing a critical threshold leading to the annihilation of one of the formerly stable states (bifurcation-induced tipping) or by a sufficiently large state perturbation (shock-induced tipping) \citep{Ashwin2012tipTypes, Alkhayuon2019AMOCRateTipping}. Rate-induced tipping is another pathway, though it will not be relevant for this work. In the bifurcation-theory framework for AMOC tipping, the transition to the off state is often conceptualised as a saddle-node bifurcation \citep{VanWesten2025SaddleNodeAMOC}. A shock-induced transition, on the other hand, can be understood as an exit of the system state from the basin of attraction associated with the AMOC on state. In both frameworks, the AMOC is investigated as an autonomous dynamical system that may be influenced by slow parameter changes or external perturbations, such as noise. This allows for the characterisation of the system in terms of several quantities, subsumed under the term "resilience" \citep{Krakovska2024ResilienceDynamicalSystems}.

One important notion of resilience has been applied to the AMOC in the context of a bifurcation-induced transition. As the system approaches such a transition through a gradual change in the background climate, the stable on-state loses linear stability, a phenomenon known as critical slowing down (CSD) \citep{Scheffer2009EWSNature}. CSD manifests in observable statistics: near the bifurcation, the dominant eigenvalue of the linearised system approaches zero, leading to increases in variance and autocorrelation of system fluctuations \citep{Held2004CSDDetection, Dakos2008CSDClimate}. These generic signatures have motivated a large body of research on early-warning indicators of tipping \citep{Lenton2011EarlyWarningClimate, Boers2021AMOCEWS, Ditlevsen2023AMOCPrediction}.

Yet CSD-based indicators capture only the system's reaction to small perturbations around the equilibrium, known as local resilience, quantifying how the on-state changes locally as the bifurcation is approached. A complementary perspective is provided by non-local resilience concepts, which characterise the size and geometry of the basin of attraction of the on-state and therefore also capture the system's capacity to recover from larger perturbations \citep{Holling1973ResilienceStabilityEcological, Gunderson2000EcologicalResilienceTheory, Menck2013BasinStability}. A comprehensive assessment of the resilience of the AMOC on state should encompass many different notions. In this work, we consider four resilience measures with the following heuristic definitions and interpretations. The concrete implementation and calculations for each model will vary and are explained in Section \ref{sec:methods}.

\begin{enumerate}
  \item \textbf{Local resilience}: capacity of a system to continuously absorb small perturbations to its equilibium on-state, such as noise and internal variability \citep{Arnoldi2016ResilienceMathComparison}.
  \item \textbf{Convergence time}: duration trajectories take to return to the on-state equilibrium after a large perturbation \citep{Neubert1997AlternativeResilience}.
  \item \textbf{Basin volume}: the volume fraction of a reference region of initial model conditions that converges to the on-state \citep{Menck2013BasinStability}.
  \item \textbf{Minimal critical shock}: the smallest instantaneous perturbation that moves the system out of the basin of attraction of the on-state \citep{Klinshov2015DistanceThresholdNumerical, Halekotte2020MinimalFatalShocks}.
\end{enumerate}

In natural systems, controlled perturbation experiments are rarely feasible, making direct estimation of non-local resilience quantities from observations difficult. This is especially true for the AMOC, and numerical models, therefore, play an essential role. However, climate models span an enormous range of complexity, from analytically tractable conceptual models to full Earth system models (ESMs). It is not obvious whether resilience measures computed in simple models translate to their more complex counterparts, which include many additional feedbacks and degrees of freedom \citep{Stouffer2006InvestigatingCausesAMOC, Schneider2017EarthSystemModeling}. If measures disagree across the hierarchy, it could mean that parts of the hierarchy do not expose them well. The extrapolation of findings from simple models to more complex ones, or to the real-world system, is then harder to justify. Several studies have examined individual resilience concepts in the context of the AMOC at specific levels of model complexity \citep{Rahmstorf2005AMOCHystersis, Cimatoribus2014AMOCBoxModel, Boers2021AMOCEWS, Ditlevsen2023AMOCPrediction, Jacques-Dumas2024RareEventResilience}, but a systematic, multi-measure comparison across a model hierarchy is lacking.

Our analysis combines four modelling studies that span many orders of magnitude of model complexity: (i) a conceptual 3-box salinity model \citep{Wood20195BoxModel, Alkhayuon2019AMOCRateTipping}, (ii) a latitude-depth Boussinesq partial differential equation model \citep{Dijkstra1997AMOCPDE, Soons2025MostLikelyNoiseinduced, Esclapez2026ConstructingEfficientScore}, (iii) the Earth system model of intermediate complexity CLIMBER-X \citep{Willeit2022EarthSystemModel}, and (iv) the coupled general circulation model PlaSim-LSG \citep{Borner2026GlobalStabilityAtlantic}. Section~\ref{sec:methods} describes each model and the concrete resilience measures computed for it. Section~\ref{sec:results} reports the individual model's results and Section~\ref{sec:discussion} synthesises them across the hierarchy. In Section~\ref{sec:discussion}, we discuss implications for research on the AMOC and on early-warning detection of potential tipping.

\section{Methods}
\label{sec:methods}

We compute a common set of resilience measures for each model while respecting the different experimental possibilities at different levels of the hierarchy. More concretely, we analyse trajectories from perturbed initial conditions. The reduced state space in which perturbations and convergence are investigated is defined by North Atlantic and Tropical Atlantic salinity values. Salt is conserved by applying a compensating perturbation in the remaining parts of the world ocean. Figure \ref{fig:perturbations-overview} illustrates the model-specific perturbation and diagnostic regions, and their implementation is explained in the respective subsections below. Table \ref{tab:overview} gives an overview of the concrete meaning and estimation routine for each resilience measure across the models. Initial conditions in the 3-box, Boussinesq, and CLIMBER-X experiments are probed within rectangles bounded by $\pm2\psu{}$ perturbations relative to their model-specific reference configurations. This is done by drawing a large number of uniform samples from this region (3-box model) or by evenly sampling a grid of initial conditions (Boussinesq model and CLIMBER-X). PlaSim-LSG is an exception here, since we analyse existing edge-track and equilibrium runs from \cite{Borner2026GlobalStabilityAtlantic} using two salinity readouts in the North and Tropical Atlantic.

We repeat the resilience analysis across a global-warming forcing coordinate, as specified for each model below. This allows us to compare directional trends of resilience, which are more meaningful across models than absolute values expressed in model-specific units. The resilience of the AMOC on-state is the primary focus throughout, though all analyses could also be carried out for other equilibria when present.

\begin{figure}[t]
  \centering
  \includegraphics[width=\textwidth]{perturbations_overview.png}
  \caption{\textbf{Perturbation and read-out geometry across the model hierarchy.}
  This overview compares the geographic salinity regions used as perturbation coordinates (shaded regions) in the Boussinesq and CLIMBER-X experiments and as salinity readouts (dashed boxes) in the Boussinesq, CLIMBER-X, and PlaSim-LSG analysis. The box geometry associated with the 3-box model is approximately reproduced from the model conception in \cite{Wood20195BoxModel}.}
  \label{fig:perturbations-overview}
\end{figure}

\begin{sidewaystable}[p]
  \centering
  \small
  \setlength{\tabcolsep}{4pt}
  \caption{\textbf{Overview of resilience-measure implementations across the model hierarchy.}}
  \label{tab:overview}

  \begin{tabular}{p{2.3cm} p{4.5cm} p{4.5cm} p{4.5cm} p{4.5cm}}
    \toprule
    & \textbf{3-Box model} & \textbf{Boussinesq} & \textbf{CLIMBER-X} & \textbf{PlaSim-LSG} \\
    \midrule
    AMOC strength
      & Overturning flux $q$ at the on-state fixed point (see SI).
      & Streamfunction mean over a fixed northern interior window of the on-state equilibrium.
      & Maximum of the Atlantic overturning streamfunction at $26°\text{N}$; trailing 500-year mean of the equilibrium run.
      & Maximum of the Atlantic overturning streamfunction between $46°$ and $66°\text{N}$; mean over the on-state equilibrium run.\\[6pt]
    Local \,\,\,\,\,\,\,\,\,\,\,\,\,\,\,\, resilience
      & $-\mathrm{Re}(\lambda_{\max})$: negative real part of the least stable eigenvalue of the Jacobian at the on-state fixed point, evaluated via forward-mode automatic differentiation \citep{Khan2015VectorForwardMode}.
      & $-\mathrm{Re}(\lambda_{\max})$: negative real part of the least stable eigenvalue fitted from a $2\times2$ discrete linear map via isotropic salinity perturbations.
      & $-\mathrm{Re}(\lambda_{\max})$: negative real part of the least stable eigenvalue fitted from a $2\times2$ discrete linear map via isotropic salinity perturbations.
      & Inverse of the long axis of the on-state fluctuation cloud ellipse.\\[6pt]
    Convergence time
      & Mean return time of on-state trajectories among $10\,000$ uniform random initial conditions in a $\pm 2\psu{}$ box; convergence at $\varepsilon=0.05\psu{}$ proximity.
      & Mean return time of valid on-state trajectories on a $20\times20$ grid over $\pm 2\psu{}$; convergence at $\varepsilon=0.05\psu{}$ proximity.
      & Mean return time of completed trajectories on a $9\times9$ grid over $\pm 2\psu{}$; convergence defined by permanent entry into an $\varepsilon=0.05\psu{}$ ball.
      & Mean transit time from the last edge-state $4\sigma$ ellipse visit to the first on-state $4\sigma$ ellipse entry. \\[6pt]
    Basin volume
      & Fraction of $10\,000$ random initial conditions in a $\pm 2\psu{}$ box centred on the on state that converge to the on state.
      & Fraction of valid initial conditions on the $20\times20$ perturbation grid converging to the on state.
      & Fraction of the $9\times9$ grid initial conditions over $\pm 2\psu{}$ that converge back to the reference equilibrium (permanent entry into an $\varepsilon=0.05\psu{}$ ball).
      & Not estimated because the available edge tracks do not thoroughly sample a reference region of state space. \\[6pt]
    Minimal \,\,\,\,critical shock
      & Minimum salinity-plane distance from the on state to a sampled initial condition not converging to the on state.
      & Minimum salinity-plane distance from the on state to a sampled initial condition not converging to the on state.
      & Minimum salinity-plane distance from the on-state to a sampled initial condition not converging to the on state.
      & Gap between the on-state and edge-state $4\sigma$ Gaussian ellipses along the line joining their centres in the deep-box salinity plane. \\
    \bottomrule
  \end{tabular}
\end{sidewaystable}



\subsection{3-Box salinity model}
\label{sec:methods:boxmodel}

We adopt the 5-box salinity model introduced by \citet{Wood20195BoxModel} and further reduced to 3 boxes by \citet{Alkhayuon2019AMOCRateTipping}. The original model partitions the Atlantic--Indo-Pacific basin into five boxes: North Atlantic ($N$), Tropical Atlantic ($T$), Southern Ocean ($S$), Indo-Pacific ($IP$), and a Deep box ($B$). In the implemented 3-box reduction, the North and Tropical Atlantic salinity anomalies, $S_N$ and $S_T$, are the two dynamical state variables. The slow variables of Southern Ocean and Deep-box salinities are prescribed as constants by the parameter vector, while Indo-Pacific salinity is derived at each time from salt conservation. The exact model equations and parameters are given in the Supplementary Information.

\citet{Wood20195BoxModel} calibrated two distinct parameter sets for the model: one for a pre-industrial climate ($p_{1\times}$) and one for a doubled-\CO{} climate ($p_{2\times}$) in the FAMOUS Earth System model. We construct a continuum of models through the full parameter vector as $p(r)=p_{1\times}+r(p_{2\times}-p_{1\times})$ and label it by $C(r)=280\,\text{ppm}+280\,\text{ppm}\,r$. Thus $r=0$ and $r=1$ correspond to 280 and 560 ppm, respectively, while $r>1$ denotes linear extrapolation beyond the doubled-\CO{} parameter vector. The nominal spacing is $\Delta r=0.05$, or $14\,\text{ppm}$, and we investigate the system from 280 through $896\,\text{ppm}$, with one additional point at $903\,\text{ppm}$ ($r=2.225$). We find that the on-state vanishes in a bifurcation when extrapolating the linear parameter ramp further, between 903 and $910\,\text{ppm}$.

All numerical experiments and resilience calculations are performed using the Julia package Attractors.jl \citep{Datseris2023FrameworkGlobalStability} within the DynamicalSystems.jl ecosystem. This computational framework enables efficient integration of large initial-condition ensembles and simultaneous tracking of multiple resilience measures as a parameter changes \citep{Morr2026ComputingResilienceMeasures}. At each continuation step, we probe the two-dimensional state space with $10\,000$ initial conditions sampled uniformly from a $\pm 2\psu{}$ rectangle centred on the on-state attractor. From these, we derive basin volume as the fraction of trajectories converging to the on state, the minimal critical shock as the minimum distance of the on state to an initial condition outside the basin, and convergence time as the mean time to reach a $0.05\psu$ neighbourhood. Local resilience is computed from the numerically derived Jacobian of the model's right-hand side at the on-state position. The resilience is measured as $-\mathrm{Re}(\lambda_\mathrm{max})$ where $\lambda_\mathrm{max}$ is the eigenvalue with the largest real part. In Figure \ref{fig:box-paper}, we map out the local basin geometry on a regular grid and show several example trajectories in terms of AMOC strength and phase-space evolution.


\begin{figure}[t]
  \centering
  \includegraphics[width=\textwidth]{amocbox_paper.png}
  \caption{\textbf{Three-box model response to salinity perturbations.}
  Top panels (a, b) show the evolution of the AMOC strength for four representative initial conditions, with dashed horizontal lines marking the on- and off-state equilibria. Bottom panels (c, d) show the same trajectories in the two-dimensional salinity state space (the two state variables), overlaid on the sampled basins of attraction; circles represent 0.05\psu{} neighbourhoods around the on- and off-state fixed points (blue and orange, respectively). The left and right columns correspond to 280 and $448\,\text{ppm}$ model specifications, respectively.}
  \label{fig:box-paper}
\end{figure}

\subsection{Boussinesq latitude-depth model}
\label{sec:methods:boussinesq}

The second model in our hierarchy is a latitude-depth Boussinesq model of the AMOC. The model evolves temperature, salinity, streamfunction, and vorticity on a $21\times41$ grid of discretised latitude and depth points. The one-dimensional sea surface is forced by cosine profiles of temperature and salinity flux, with lower temperatures and lower net freshwater input at high latitudes. This model is much less expensive than global circulation models but already contains spatially structured salinity fields and overturning, making it a useful intermediate step between the conceptual box model and the three-dimensional climate models.

We introduce a single external forcing parameter that mimics Arctic amplification. Increasing it flattens the equator-to-pole surface temperature profile (Arctic amplification) and enhances the freshwater forcing profile (intensification of the hydrological cycle). For the exact mechanism of this parameter, its mapping to \CO{}, and the exact model equations, consult the Supplementary Information. The perturbation experiments are performed across seven forcing values corresponding to \CO{} concentrations from $280\,\text{ppm}$ to approximately $643\,\text{ppm}$. For each configuration, the model is first spun up twice under weak hemispheric salinity-forcing asymmetries to ensure convergence to the on- and off-states, respectively. Under the symmetric salinity forcing, for which an active and a reversed overturning state can coexist over part of the forcing range, we obtain the final AMOC on and off reference states by evolving the asymmetric-forcing equilibria until convergence. Based on the reference on-state, we initialise perturbed salinity trajectories. We then classify each perturbed run by the reference state it settles to, which in turn enables the basin-volume, minimal critical shock, and convergence-time measures.

Initial-condition perturbations are applied to grid-aligned salinity boxes, and we define an affine conversion of the non-dimensional salinity units of the Boussinesq model to \psu{} (for details see Supplementary Information). The North Atlantic box covers $40.5°$--$90°\text{N}$, the Tropical Atlantic box covers $31.5°\text{S}$--$40.5°\text{N}$, and the Southern compensation box covers the domain south of $31.5°\text{S}$. All three extend approximately through the upper third of the water column. The perturbation fields are cosine-tapered to avoid sharp salinity gradients. Perturbations are applied as salinity anomalies $(\Delta S_N,\Delta S_T)$ on a $20\times20$ grid ($400$ initial conditions), with both amplitudes spanning $\pm2\,\psu{}$ relative to the asymmetric-forcing on-state equilibrium, which lies within $0.1\,\psu{}$ of the final on-state. A compensating perturbation in the Southern box conserves total domain salinity.

For the basin-mapping ensemble, only the salinity field of the on-state equilibrium is carried over into the perturbed initial conditions. The temperature, vorticity, and streamfunction fields are set to zero and left to re-adjust to the perturbed salinity under the model dynamics. This isolates the salinity perturbation as the driver of the response. To estimate local linear stability, we run a second, smaller ensemble of eight experiments per forcing level. Each run has a salinity perturbation in the two-dimensional (North and Tropical Atlantic) salinity space. The eight perturbations are evenly arranged on a circle of radius $\varepsilon = 0.5\,\psu{}$ around the on-state equilibrium. Unlike the basin-mapping ensemble, these perturbations are applied to the complete equilibrium state, retaining all temperature, vorticity, and streamfunction fields, so that the local recovery of the full system is reflected in the two-dimensional salinity projection. All trajectories are run under neutral symmetric salinity forcing for approximately $500$ model years. After discarding an initial ${\approx}50$-year transient, a $2\times 2$ discrete linear map $M$ is fitted by least squares to trajectory pairs $(x_t,\, x_{t+\Delta t})$ sampled every ${\approx}20$ years, where $x = (\delta S_N,\, \delta S_T)$ denotes the box-salinity deviation from equilibrium. The continuous-time generator $A = \log(M)/\Delta t$ of the fitted map $M$ yields eigenvalues whose least stable real part gives the local resilience $-\mathrm{Re}(\lambda_{\max})$.

Time metrics are reported on a year-proxy scale, set to 100 times the arbitrary model time unit, to facilitate comparability across illustrations. In Figure \ref{fig:boussinesq-paper}, we plot four exemplary trajectories in terms of the two box salinities and the AMOC strength at $280\,\text{ppm}$ and approximately $488\,\text{ppm}$.


\begin{figure}[t]
  \centering
  \includegraphics[width=\textwidth]{boussinesq_paper.png}
  \caption{\textbf{Boussinesq model response to salinity perturbations.}
  Top panels (a, b) show the evolution of the AMOC strength for four representative initial conditions (the first 50 years of initial adjustment are omitted), with dashed horizontal lines marking the on- and off-state equilibria. Bottom panels (c, d) show the same trajectories in North Atlantic and Tropical Atlantic box-mean salinity, overlaid on the sampled basins of attraction. Circles represent 0.05\psu{} neighbourhoods around the on- and off-state fixed points (blue and orange, respectively). Trajectories end once they enter one of these neighbourhoods, which can happen before the AMOC strength has fully adjusted. The left and right columns correspond to 280 and $488\,\text{ppm}$ model specifications, respectively.}
  \label{fig:boussinesq-paper}
\end{figure}


\subsection{CLIMBER-X}
\label{sec:methods:climberx}

CLIMBER-X \citep{Willeit2022EarthSystemModel} is a fast Earth system model of intermediate complexity (EMIC) that couples a 3-D frictional-geostrophic ocean model, a 2.5-D statistical-dynamical atmosphere, a dynamic-thermodynamic sea-ice model, and land-surface components. The ocean component resolves the full three-dimensional salinity and temperature fields, including the main overturning circulation cells, at a resolution of approximately $5°\times 5°$ with 23 vertical levels. CLIMBER-X is computationally efficient enough to permit multi-millennial integrations, making it well-suited for exploring the state space through moderately sized perturbation ensembles.

We first equilibrate CLIMBER-X at 15 prescribed atmospheric \CO{} concentrations between 285 and $600\,\text{ppm}$. Each equilibrium run is integrated for up to $20\,000$ model years from the default pre-industrial restart. After equilibration, the trailing $500$ years of the annual ocean time series are used to store reference values for AMOC strength at $26°\text{N}$ and for the salinity diagnostics in the geographical boxes illustrated in Fig.~\ref{fig:perturbations-overview}.

The North Atlantic box comprises Atlantic-basin grid cells from $35°$--$80°\text{N}$ and the Tropical Atlantic box comprises Atlantic-basin cells from $35°\text{S}$--$35°\text{N}$. Both boxes cover the upper $105\,\text{m}$, corresponding to the top four ocean layers. The North and Tropical box perturbations use a two-grid-cell latitudinal taper and a two-layer vertical cosine taper. The tapered fields are normalised so that the volume-weighted mean anomaly over the hard target box remains exactly equal to the prescribed salinity perturbation. The salt added to the two Atlantic boxes is compensated uniformly over the rest of the world ocean within the same upper $105\,\text{m}$, so that the global ocean salt content in the restart is conserved. 

For our experiments, we apply instantaneous salinity perturbations to these boxes directly to the ocean restart as salinity anomalies in \psu{}, creating an ensemble of modified initial conditions. The North Atlantic and Tropical Atlantic perturbations each take the values $-2,-1.5,\ldots,2\psu{}$, giving a $9\times9$ grid of 81 experiments per \CO{} level. Each perturbed initial condition is integrated forward for $10\,000$ model years. This model is known to exhibit an AMOC off-state under high \CO{} concentrations when large constant freshwater forcing in the North Atlantic is applied. In the present setup of pure initial salinity perturbations, the off-state is not reached. Instead, at higher \CO{} a second, distinct attractor with a somewhat stronger AMOC emerges. This occurs when sea ice in the Northern Hemisphere retreats and the deep-water convection driving the AMOC expands into areas farther north \citep{Willeit2024GeneralizedStabilityLandscape}. We classify each grid trajectory by its final state using the same criterion as for the Boussinesq model: a run is assigned to the on-state if its final salinities lie within an $\varepsilon=0.05\,\psu{}$ ball of the unperturbed reference state, and to the ``strong'' state otherwise.

A further \CO{}-dependence concerns the long equilibration run itself: the state on which it settles, and which serves as the unperturbed reference, changes branch with forcing. Up to $330\,\text{ppm}$ it lies on the reference (``modern'') branch, whereas from $345\,\text{ppm}$ onward it lies on the stronger-AMOC (``strong'') branch. The resilience measures obtained from the perturbation experiments are therefore associated with the modern state at $285$--$330\,\text{ppm}$ and with the strong state at $345$--$600\,\text{ppm}$. The basin volume reported for CLIMBER-X (Table~\ref{tab:overview}) is the share of the 81 grid trajectories that return to this reference equilibrium. It equals one wherever all perturbations recover and falls below one only near the modern-branch transition (at $315$ and $330\,\text{ppm}$), where a fraction of perturbations settle on the strong state instead. From the same grid, we obtain the minimal critical shock as the smallest salinity-plane distance from the reference equilibrium to a perturbed initial condition that does not return to its basin. It is defined only where at least one perturbation escapes, i.e.\ at $315$ and $330\,\text{ppm}$, since at all other \CO{} levels the entire $\pm2\psu{}$ grid recovers. Figure \ref{fig:climberx-paper} shows four example trajectories starting from perturbed salinity initial conditions in the specified boxes, similar to the two figures above. For both \CO{} values, the equilibrium states are the modern state.

To estimate local linear stability, we run a second, smaller ensemble of $9$ experiments per \CO{} level. Eight runs apply isotropic salinity perturbations of magnitude $\varepsilon = 0.1\,\psu{}$ again arranged on a circle in the two-dimensional (NA, Tropical) shallow-box salinity space, plus one unperturbed reference run. Each run is integrated for $500$ model years. An equilibrium reference is obtained from the mean of the unperturbed run over years 10--150. A $2\times2$ discrete linear map is fitted by least squares to annual trajectory pairs $(x_t,x_{t+1})$ in the same fitting window. The continuous-time generator $A=\log(M)/(1\,\text{yr})$ yields eigenvalues whose least stable real part gives the local resilience $-\mathrm{Re}(\lambda_{\max})$.


\begin{figure}[t]
  \centering
  \includegraphics[width=\textwidth]{climberx_paper.png}
  \caption{\textbf{CLIMBER-X model response to salinity perturbations.}
  Top panels (a, b) show the AMOC strength at $26°\text{N}$ along four example trajectories from perturbed salinity initial conditions. Bottom panels (c, d) show the same trajectories in the two-dimensional Atlantic-box salinity plane, with axes spanning exactly the $\pm2\,\psu{}$ perturbation window. The shaded background maps the basin of attraction of every initial condition on the full $9\times9$ grid: cells whose trajectory returns to the reference on-state (final salinity within $\varepsilon=0.05\,\psu{}$) are blue, and cells that settle on the distinct stronger-AMOC (``strong'') state are purple. Circles mark the $\varepsilon=0.05\,\psu{}$ neighbourhoods of the on-state (blue, at the unperturbed equilibrium) and of the strong state (purple, at the mean endpoint of the non-converged trajectories). The two columns show prescribed atmospheric \CO{} concentrations of $285\,\text{ppm}$ and $330\,\text{ppm}$: at $285\,\text{ppm}$ all 81 initial conditions return to the on-state, whereas at $330\,\text{ppm}$ 46 of the 81 converge to the strong state, which lies only ${\approx}0.08\,\psu{}$ from the on-state but exhibits an AMOC about $0.4\,\Sv{}$ stronger. Dashed horizontal lines in (a, b) mark the equilibrium AMOC strength of the on-state (blue) and, in (b), of the strong state (purple).}
  \label{fig:climberx-paper}
\end{figure}

\subsection{PlaSim-LSG}
\label{sec:methods:plasim}

The last component of the hierarchy is the coupled climate model of intermediate complexity PlaSim-LSG, comprising a spectral primitive-equation atmosphere (T21, $\approx 5.6°\times 5.6°$, 10 vertical levels) and the Large Scale Geostrophic (LSG) ocean \citep{Maier-Reimer1993MeanCirculationHamburg}, a full 3-D ocean circulation model at $3.5°\times 3.5°$ resolution with 22 vertical levels, coupled to a thermodynamic sea-ice scheme and a hydrological cycle.  With approximately $10^5$ degrees of freedom and a throughput of roughly 700 model years per CPU day, PlaSim-LSG is considerably more complex than CLIMBER-X. In particular, it resolves atmospheric variability more faithfully and thereby yields seemingly chaotic ocean dynamics.

At $285\,\text{ppm}$ (pre-industrial) and $360\,\text{ppm}$ (contemporary) \CO{}, PlaSim-LSG exhibits a bistable AMOC with an on-state of $\approx 15\Sv{}$ and an off-state (deep convection in Subpolar Gyre ceases) of $\approx 1-2\Sv{}$, both confirmed by long unforced simulations. \citet{Borner2026GlobalStabilityAtlantic} applied an edge-tracking algorithm \citep{Lucarini2017EdgeTracking} to locate the edge state (Melancholia state) of the AMOC in PlaSim-LSG at $285\,\text{ppm}$ and $360\,\text{ppm}$. This is a quasi-stable equilibrium on the complex basin boundary between the on and off states. The algorithm iteratively bisects pairs of trajectories converging to opposite attractors and tracks the resulting pseudotrajectory along the basin boundary until it converges to the edge state. The edge state is a chaotic saddle characterised by persistent centennial AMOC oscillations ($\approx 120$\,yr period, amplitude $3$–$10\Sv{}$), driven by atmosphere–ice–ocean interactions in the North Atlantic. A complete description of the model configuration, its AMOC bistability, and the edge-tracking procedure is given in \citet{Borner2026GlobalStabilityAtlantic}, to which we refer the reader for full methodological details.

From the final iterations of edge tracking at each \CO{} level, we extract pairs of edge tracks: one converging to the on-state and one converging to the off-state, both initialised from near-boundary conditions that closely shadow the edge state before diverging. These trajectories provide a natural set of initial conditions from which resilience-relevant transient behaviour can be quantified: they start close to the basin boundary and thus represent challenging recovery paths for the on-state. In total, we analyse 37 paired edge-track files at $285\,\text{ppm}$ and 38 at $360\,\text{ppm}$. Dedicated long equilibrium simulations of the AMOC-on, AMOC-off, and edge states at each \CO{} level provide the attractor clouds against which these transients are evaluated.

We reduce the high-dimensional PlaSim-LSG salinity output to the salinity of two geographical boxes. The North Atlantic readout is the mean salinity from $35°$--$80°\text{N}$ and 0--$1000\,\text{m}$ depth, and the Tropical Atlantic readout covers $35°\text{S}$--$35°\text{N}$ and 0--$500\,\text{m}$ depth. Dedicated on-, off-, and edge-state equilibria trajectories determine the means associated with these states in this two-dimensional salinity plane. The covariance matrices of the scattered trajectory points are computed to obtain an approximate Gaussian distribution of data for each state. Convergence time is measured from the last visit to the edge-state $4\sigma$ Gaussian ellipse to the first entry into the target-attractor ellipse, and the critical-shock proxy is the gap between the edge-state and on-state ellipses, measured along the line joining their centres. Local resilience is quantified as the inverse width of the on-state fluctuation cloud, $[2000\sqrt{\lambda_{\max}(C)}]^{-1}$, where $\lambda_{\max}(C)$ is the largest eigenvalue of the on-state equilibrium covariance $C$. A broader cloud, i.e.\ larger salinity fluctuations around the on-state, corresponds to lower local resilience since the system appears less capable of absorbing internal variability.

\begin{figure}[t]
  \centering
  \includegraphics[width=\textwidth]{plasim_paper.png}
  \caption{\textbf{PlaSim-LSG edge-state trajectories and covariance ellipses of attractors.}
  Top panels (a, b) show the AMOC strength readout along converged edge-track trajectories. Bottom panels (c, d) show the same transients in the North and Tropical Atlantic salinity coordinates. Ellipses denote the fitted $4\sigma$ Gaussian regions to equilibrium runs of the three states (on: blue, off: orange, edge: black), used to define convergence times and edge-to-attractor distances. Trajectories are coloured by the state they converge to. The two columns show the experiment results at prescribed atmospheric \CO{} concentrations of $285\,\text{ppm}$ and $360\,\text{ppm}$, respectively.}
  \label{fig:plasim-paper}
\end{figure}

\section{Results}
\label{sec:results}

The trajectories from the described perturbation experiments are synthesised to give resilience estimates according to all measures listed in Table~\ref{tab:overview}. The evolution of the resilience measures is shown in Figure~\ref{fig:synthesis}, which forms the basis for the subsequent comparison and differentiation of the results. The models differ substantially in the dynamical characteristics of the AMOC, including its stationarity, multistability, and bifurcations. This, in itself, complicates the interpretation of our results on AMOC resilience within a common framework.

The models do not agree on a consistent impact of atmospheric \CO{} on the equilibrium AMOC strength. The 3-box model exhibits AMOC weakening, which is consistent with behaviour exhibited by the FAMOUS Earth System Model \citep{Bouttes2015NonlinearityOceanHeat}, the calibration target for the 3-box model \citep{Alkhayuon2019AMOCRateTipping}. The vanishing of the AMOC on-state in a bifurcation at a threshold \CO{} concentration has, on the other hand, not been observed in FAMOUS. In the Boussinesq model, CLIMBER-X, and PlaSim-LSG, AMOC strength increases as the \CO{} concentration rises. For the latter two models, it has been conjectured that this is due to a changing net surface freshwater flux, particularly in the Tropical Atlantic \citep{Willeit2024GeneralizedStabilityLandscape}. More evaporation implies more saline water reaching the North Atlantic. Abstracting this mechanism could provide insight into the same phenomenon in the Boussinesq model. Because the surface freshwater forcing weakens in the central (tropical) part of the model relative to the right (northern) part, the salt-advection feedback might weaken as a response to the implicit \CO{} forcing. None of the three models exhibits a vanishing AMOC on-state in a bifurcation.

The existence and characteristics of an AMOC off-state are also not consistent across the model hierarchy. The 3-box and Boussinesq models feature stable states with reversed AMOC flow. That is, the salinity gradient between the North and Tropical Atlantic reverses in both models, and in the Boussinesq model, the actual fluid overturning changes direction. It is unclear whether this is physically meaningful behaviour, i.e. whether any real-world AMOC collapse would lead to a strong reversed flow. However, as we regard alternative stable states primarily as differentiations from the contemporary AMOC state, their exact characteristics are of little concern. The PlaSim-LSG model exhibits an AMOC off-state characterised by a breakdown of deep-water convection in the North Atlantic. The CLIMBER-X model exhibits a weak AMOC state only at \CO{} concentrations lower than preindustrial conditions. This is due to more extensive North Atlantic sea ice blocking deep-water convection in the Labrador and Irminger Seas. The only relevant alternative state for our experiments is a stronger AMOC state at large \CO{} concentrations, reached when diminished Northern sea ice allows for additional deep-water convection in the Kara Sea and Nansen Basin \citep{Willeit2024GeneralizedStabilityLandscape}. It is important to consider these differences, especially when interpreting the results of basin volume and minimal critical shock. The two measures should only be interpreted as resilience against AMOC collapse when the alternative state is indeed associated with an AMOC weakening.

Local resilience of the on-state describes the restoring force that returns small perturbations back to the equilibrium. It can be interpreted as an upper bound on the reciprocal of the half-life of any small perturbation. Critical slowing-down theory predicts a decrease in local resilience as one approaches a bifurcation tipping point. This is indeed what we see for the 3-box model. Local resilience of the Boussinesq and CLIMBER-X AMOC also seems to decrease with increasing \CO{}, though the decrease is less pronounced and accompanied by fluctuations in the latter case. In the PlaSim-LSG model, the trajectory cloud forming the quasistationary on-state equilibrium slightly decreases in size, which is consistent with a larger local Lyapunov exponent and, thereby, stronger local resilience.

The convergence time aggregates show consistent trends as \CO{} forcing increases. In all models, the mean time needed for a probed perturbation to converge back to the on-state increases. This can be interpreted as a loss of resilience, as the AMOC is outside its usual operating regime for longer. One model-specific behaviour is worth emphasising in this context. The transient trajectories of some perturbations in the CLIMBER-X model at 315ppm show long occupancy of the strong AMOC state before eventually converging to the modern state (see SI Figure S1). It is possible that all perturbations marked as converging to the strong state in the 10,000 years of trajectory integration would eventually return to the modern state. Under the 330ppm configuration, these long transients remain in the strong state basin for the entire 10ky, reducing the estimated basin volume of the modern state and the mean convergence time. Whether this, too, constitutes spurious meta-stability cannot be determined with these experiments. This highlights the limitations in the numerical computation of resilience measures and the need to accurately define a finite-time framework for their interpretation. Throughout this study, convergence is never determined analytically but obtained via a finite-time cutoff evaluation against a neighbourhood convergence criterion \citep{Morr2026ComputingResilienceMeasures}.

A less intuitive result emerges from the basin volume computations. The share of initial perturbations in the region spanned by $\pm2$\psu{} in Northern and Tropical salinity that converge to the on-state rises in the 3-box and Boussinesq model when \CO{} increases. This is surprising insofar as it indicates a strengthening of the AMOC's resilience in this interpretation. Only in the immediate proximity to the bifurcation threshold of the 3-box model do we see a sharp decrease in basin volume. The basin structure of the CLIMBER-X model is almost trivial. For most considered \CO{} concentrations, all probed perturbations converge to either the modern AMOC on-state (300ppm and below) or the emerging strong AMOC state (345ppm onwards). Only at 315 ppm and 330 ppm do perturbations of the modern state lead to a transient convergence to the strong state, splitting the basin.

Finally, the minimal critical shock shows a consistent increase in the 3-box, Boussinesq, and PlaSim-LSG models. In CLIMBER-X, it is only defined at 315 and $330\,\text{ppm}$, where it decreases as the strong state takes over a growing part of the probed region. The increase in the other three models would mean an increase in resilience, as larger perturbations are needed to push the AMOC to divergence away from the on-state.


\begin{figure}[t]
  \centering
  \includegraphics[width=\textwidth,height=0.82\textheight,keepaspectratio]{synthesis_figure.png}
  \caption{\textbf{Cross-model synthesis of AMOC resilience measures.}
  Equilibrium AMOC strength (a) and Tropical minus North Atlantic salinity gradient (b). Resilience quantifiers are plotted against the \CO{} coordinate for each model, where available (c--f). The model-specific definitions and implementation of each resilience quantifier are given in Table \ref{tab:overview}. These results provide the basis for assessing which resilience signals are coherent across the model hierarchy.}
  \label{fig:synthesis}
\end{figure}

\section{Discussion}
\label{sec:discussion}

The picture emerging from our multi-model, multi-measure comparison of AMOC resilience is ambiguous. It is confounded by challenges in conceptual compatibility and practical differences in dynamical characteristics. Some measures exhibit internal inconsistencies in resilience trends across models. Other measures show a consistent increase in resilience, which seems incompatible with a consistent decrease in resilience on another measure. The notion of a real-world AMOC at risk of bifurcation tipping due to anthropogenic influences suggests a decrease in all resilience measures discussed near the bifurcation point \citep{Morr2026ComputingResilienceMeasures}. The simple bifurcating 3-box model does produce a decrease in basin volume and minimal critical shock, but only as it enters the direct proximity of the bifurcation, and the on-state basin contracts to a smaller size than the $\pm2$\psu{} salinity region.

The moderate decrease in local resilience observed in the Boussinesq and CLIMBER-X models reinforces the notion that critical slowing down should be expected prior to a bifurcation-mediated AMOC collapse. However, these quantities are obtained in equilibrium, and the real-world AMOC is likely still adapting to fast anthropogenic influences. The direct implications for the applicability of CSD-based early warning signals, therefore, remain limited. In addition, the PlaSim-LSG result suggests a strengthening of local resilience through a reduction in the stationary variance of the quasi-equilibrium state. The consistent increase in non-local convergence time could be interpreted as indicating that the change in the local restoring force generalises to larger regions of the basins of attraction. It lends further credence to the notion of critical slowing down, even as the local radius of perturbations is extended.

Basin volume, as we have computed it here, provides insight into how the on-state provenance evolves under the chosen perturbation regime. For the Boussinesq model, the increase is moderate, as in fact only 31 out of 400 perturbation runs switch their basin membership between 280 and $643\,\text{ppm}$, all of them from the off- to the on-state basin. The contraction of the modern state basin in the CLIMBER-X model reveals novel insights into Earth System Model dynamics. To our knowledge, this is the first instance of a model of this size exhibiting bistability with respect to changes in the initial salinity configuration. Typically, modelling studies employ persistent hosing influences to permanently alter surface salinity and facilitate an AMOC collapse. The possibility that pulse perturbations could similarly lead to the convergence to an alternative stable state, albeit a stronger one, should inspire future research.

The minimal critical shock provides context for the results on basin volume, since the two are mechanistically related. It reveals that the increase in basin volume is at least in part mediated by an expansion of the basin boundary at the point where it is closest to the on-state. This point is a natural candidate for the so-called instanton path to pass through \citep{Grafke2019NumericalComputationRare, Soons2025MostLikelyNoiseinduced}. This is the most likely path in which a random transition occurs when forced by small noise. It is also related to saddle points or edge states, which are unstable equilibrium points on the boundary \citep{Lucarini2017EdgeTracking, Lohmann2025EdgeStatesEWS}. The result of the PlaSim-LSG model is consistent with the interpretation of a distancing saddle point on the basin boundary of the 3-box and Boussinesq model. This is because, in this case, we are measuring the distance between the on-state and the edge state, which was directly approximated using an edge-tracking algorithm \citep{Borner2026GlobalStabilityAtlantic}. The consistent increase in minimal critical shock should be studied in a broader selection of models going forward, as it reveals important information about the most likely transition path in an AMOC collapse.

Our comparative study of AMOC resilience does not offer concrete physical insight into AMOC dynamics. However, the discrepancies and consistent behaviours found across measures and models inspire a more careful selection of these measures and models. In principle, the evolution of resilience measures could provide valuable insight into the mechanistic changes the AMOC undergoes under climate change. It is not a priori expected that simple models exhibit the same resilience behaviour as complex models, and none may represent the true AMOC dynamics well enough to allow generalisation of our findings in any sense. Our findings offer a first insight into whether and how resilience can be consistently defined and computed across models of varying complexity. Experiments with higher-fidelity models may help resolve some of the discrepancies that arose here. If the real-world AMOC is indeed approaching a bifurcation threshold, we expect insights from the simple 3-box model to be valuable, as many current climate models struggle to reproduce this threshold \citep{Valdes2011BuiltStability, Liu2017OverlookedPossibilityCollapsed, McCarthy2023TrustAMOCProjections}. On the other hand, results from simple models should be regarded with caution and treated as strong abstractions.

We hope to inspire future research on the AMOC and other natural systems under the framework of dynamical system resilience. Its system-agnostic nature enables effective comparison of system dynamics, while questions remain about its consistent adaptability across model hierarchies.


\bibliographystyle{apalike}
\bibliography{BibAll}


\end{document}
